\documentclass{article}
\usepackage[utf8]{inputenc}
\usepackage{amsmath}
\usepackage{geometry}
\geometry{margin=2cm}
\begin{document}

\section*{Questão 119 --- ENEM 2024 --- Ciências da Natureza}

\subsection*{Interpretação da Questão e Estratégia}

O enunciado apresenta a função das zonas de deformação (\emph{crumple zones}) em carros modernos e questiona por que elas aumentam a segurança dos passageiros durante uma colisão.

\paragraph{Termos e conceitos chave}
\begin{itemize}
  \item \textbf{Considerar}: analisar o sistema descrito (carro, ocupantes e muro) isoladamente para entender o fenômeno físico.
  \item \textbf{Deformação}: alteração da forma estrutural do veículo para absorver energia do impacto.
\end{itemize}

\paragraph{O que a questão pede}
Explicar como a deformação da estrutura do carro por meio da zona de deformação contribui para a segurança dos ocupantes durante a colisão, no contexto do sistema isolado.

\paragraph{Estratégia de resolução}
Deve-se analisar a colisão considerando conservação da quantidade de movimento e transformação da energia cinética em energia de deformação. A zona deformável absorve essa energia, aumentando o tempo de impacto e assim reduzindo a força aplicada nos passageiros, o que lhes protege. Compreender a diferença entre quantidade de movimento (conservada) e energia cinética (dissipada) é essencial para responder corretamente.

\bigskip

\subsection*{Revisão Conceitual}

\begin{tabular}{|p{5cm}|p{6cm}|p{6cm}|}
\hline
\textbf{Conceito} & \textbf{Ideia-chave} & \textbf{Aplicação no problema} \\
\hline
Conservação da quantidade de movimento & Em sistemas isolados, a quantidade de movimento total é constante durante a colisão. & Sistema carro, ocupantes e muro desacelera, mas mantém a quantidade de movimento constante. \\
\hline
Energia cinética & Energia que depende da velocidade do corpo e pode se transformar em outras formas. & Energia cinética do carro é dissipada na deformação das \emph{crumple zones}. \\
\hline
Tempo de colisão e força média & Aumentar o tempo de colisão diminui a força média sobre os ocupantes (\(F = \Delta p/\Delta t\)). & A zona deformável aumenta o tempo do impacto, reduzindo a força sofrida pelos passageiros. \\
\hline
\end{tabular}

\bigskip

\paragraph{Fundamentos teóricos}
\begin{itemize}
  \item Em uma colisão, considerando-se o sistema isolado composto por carro, ocupantes e muro, a quantidade de movimento total se conserva.
  \item A energia cinética do carro é convertida em outras formas, como a energia de deformação na zona deformável.
  \item A deformação aumenta o tempo durante o qual ocorre a colisão, reduzindo a força média sobre os ocupantes.
\end{itemize}

\bigskip

\subsection*{Resolução da Questão e \emph{Pulo do Gato}}

\begin{enumerate}
  \item Identificar que a zona de deformação (\emph{crumple zone}) é projetada para deformar-se durante a colisão.
  \item Com a deformação, ocorre dissipação da energia cinética, que não é transmitida integralmente aos ocupantes, protegendo-os de forças muito intensas.
  \item O sistema isolado conserva a quantidade de movimento, mas a energia cinética é transformada em energia de deformação e outras formas.
  \item A dissipação da energia faz com que o tempo total da colisão aumente, diminuindo a força média atuante sobre os ocupantes.
  \item Conclui-se que a deformação permite absorber a energia do impacto, protegendo os passageiros, confirmando assim a alternativa \textbf{B}.
\end{enumerate}

\paragraph{Pulo do gato}
A chave para resolver a questão é entender que aumentar o tempo da colisão pela deformação reduz a força sobre os passageiros, aumentando a segurança.

\bigskip

\subsection*{Análise dos Distratores}

\begin{itemize}
  \item \textbf{A}: Errado. O \emph{crumple zone} não aciona diretamente os airbags, que funcionam por sensores distintos.
  \item \textbf{B}: Correto. A região deformável absorve parte da energia cinética, reduzindo a força transmitida aos ocupantes.
  \item \textbf{C}: Errado. A quantidade de movimento do sistema se conserva; a deformação absorve energia, não quantidade de movimento.
  \item \textbf{D}: Errado. A zona de deformação não é uma barreira rígida, mas uma parte que se deforma para absorver energia.
  \item \textbf{E}: Errado. A velocidade do centro de massa do sistema isolado é conservada; a deformação dissipa energia, mas não altera essa velocidade.
\end{itemize}

\bigskip

\subsection*{Código da Questão}
CN\_FISICA\_MULTIPLAESCOLHA\_001

\bigskip

\textit{\small Imagens do enunciado ilustram a zona deformável \emph{crumple zone}, os airbags e os cintos de segurança, ressaltando a deformação da estrutura dianteira do carro na colisão contra o muro.}

\end{document}